We report the broad measured deployment selection, retaining quality-loss, single-repeat and screening observations with their original flags. Throughput is the recorded cohort median in images/\allowbreak{}s. Energy is calibrated host-inclusive energy per completed image, averaged as the mean of per-repeat E/\allowbreak{}N and displayed in mJ/\allowbreak{}image. Dataset Top-1 accuracy or COCO mAP is reused from the selected existing evidence and shown separately. Table 3 ranks Native splits retrospectively against the fastest measured Full on the same host and uses that exact Full for both speedup and energy gain. Source generations and numerical quality/\allowbreak{}cache discrepancies are recorded per metric. Empirical performance--energy fronts maximize throughput and minimize host-inclusive energy within each model/\allowbreak{}setup, including both Full alternatives; quality is shown by color and is not a dominance objective. The 75 updated Native DeepX observations use batch-one, single-engine A4 with at most four outstanding jobs; splits use FIFO depth three and a serial TensorRT consumer through task completion. The other 375 direct observations and all 450 selected dataset-quality observations are retained. Native detector splits record a correctness oracle outside timing, whereas Generic detector splits record decoded NMS; Full paths retain their recorded host-tail or integrated normalization work. These ratios compare observed deployments with their recorded parallelism and do not isolate a causal effect of splitting at identical work.
